% =============================================================================
% Bloque para pegar en latex/informe.tex
% después de \subsection{Análisis por caso} dentro de la sección Resultados.
%
% Si no están ya en el preámbulo, añadir:
%     \usepackage{xcolor,colortbl}
%     \usepackage{booktabs}
%     \usepackage{listings}
%     \definecolor{cellgood}{RGB}{212,237,218}
% =============================================================================

\subsection{Métricas de calidad contra GT real}
\label{subsec:metricas-gt}

Para cuantificar la fidelidad geométrica del interpolado contrasté, para cada
par $(A, B)$, el contorno reconstruido a $t = 0.5$ contra el slice real
intermedio (\emph{ground truth}) del dataset BraTS. Reporto cuatro métricas:
\emph{Dice} y \emph{IoU} de solapamiento de región (mejor cuanto mayor;
calculadas rasterizando ambos polígonos en una grilla $512\!\times\!512$),
distancia de Hausdorff bidireccional entre vértices (mejor cuanto menor) y
error relativo de área $\mathrm{areaErr}=|A_{\text{interp}}-A_{\text{gt}}|/A_{\text{gt}}$.

\begin{table}[h]
  \centering
  \small
  \begin{tabular}{lcccccc}
    \toprule
    Caso & Par $(A \!\to\! B)$ & Dice & IoU & Hausdorff & AreaErr & Self-int \\
    \midrule
    00008-100 & $70 \!\to\! 72$ &
      \cellcolor{cellgood}0.855 & \cellcolor{cellgood}0.746 &
      \cellcolor{cellgood}1.265 & \cellcolor{cellgood}0.142 & No \\
    00008-101 & $70 \!\to\! 72$ &
      \cellcolor{cellgood}0.854 & \cellcolor{cellgood}0.746 &
      \cellcolor{cellgood}0.896 & \cellcolor{cellgood}0.056 & No \\
    \midrule
    \textbf{Promedio} & --- &
      \textbf{0.855} & \textbf{0.746} &
      \textbf{1.081} & \textbf{0.099} & --- \\
    \bottomrule
  \end{tabular}
  \caption{Métricas del interpolado lineal a $t=0.5$ contra el slice GT real
  intermedio. El sombreado verde indica que la celda cumple el umbral de
  calidad (Dice $\geq 0.80$, IoU $\geq 0.65$, Hausdorff $\leq 2.0$,
  AreaErr $\leq 0.20$).}
  \label{tab:metricas-gt}
\end{table}

Es importante remarcar una diferencia metodológica con el trabajo de Santiago.
Él reporta $\mathrm{Dice} \approx 0.997$ sobre el dataset completo, midiendo
contra el GT \emph{sintético} reconstruido del propio dataset (es decir,
contra la misma malla que su pipeline produce); por construcción esa
comparación tiende a saturar en 1. En cambio, mi protocolo descarta el slice
intermedio del par y mide el interpolado contra el GT \emph{real} de ese
slice, lo que castiga cualquier desviación topológica o de forma. Ambos
enfoques son válidos y complementarios, pero no son directamente comparables:
los valores $\sim\!0.85$ que reporto reflejan la dificultad intrínseca de
predecir un slice ausente a partir de sus vecinos.


\subsection{Experimento de stress sobre el resolver}
\label{subsec:stress}

Para validar el módulo \texttt{SelfIntersectionResolver} diseñé un experimento
que fuerza una correspondencia patológica entre $A$ y $B$:

\begin{enumerate}
  \item Tomo el par del caso 00008-100 (\texttt{slice\_0070.obj},
        \texttt{slice\_0072.obj}).
  \item Rotación cíclica de los vértices de $B$ por $k = n/2$ posiciones,
        donde $n$ es la cantidad de vértices de $B$. Esto desalinea por
        completo la correspondencia angular inicial y, sin compensación,
        forzaría que los segmentos del interpolado a $t=0.5$ se crucen
        entre sí.
  \item Interpolo a $t = 0.5$ con el binario y, si aparecen
        auto-intersecciones, las resuelvo con \texttt{SelfIntersectionResolver}.
\end{enumerate}

\begin{figure}[h]
  \centering
  \includegraphics[width=0.95\textwidth]{../src/interpolation-lineal/output/stress_test/2026-05-31_16-48-37/stress_comparison.png}
  \caption{Comparativa lado a lado del experimento de stress sobre
  \texttt{BraTS-GLI-00008-100}. Izquierda: corrida normal con $B$ original.
  Derecha: $B$ rotado en $k = n/2 = 12$ vértices. En ambos casos el
  contorno verde es el interpolado a $t=0.5$. La búsqueda de rotación
  óptima del binario, seguida del resolver, recupera una topología
  válida incluso bajo la rotación adversarial.}
  \label{fig:stress}
\end{figure}

El experimento demuestra que el pipeline completo
(rotación óptima $+$ resampling $+$ interpolación $+$ resolver) es robusto
frente a perturbaciones del orden cíclico de los vértices de entrada, que es
exactamente el tipo de inconsistencia esperable cuando los contornos vienen
de slices independientes sin un punto de inicio canónico.


\subsection{Salida para reconstrucción 3D}
\label{subsec:multi-t}

El binario escribe siempre los vértices con $z = 0$, lo cual es suficiente
para inspeccionar el resultado en 2D pero impide apilar varios contornos
interpolados en un volumen tridimensional para reconstrucción posterior
(p.ej.\ Poisson o marching cubes en VTK/ParaView). Para cerrar ese hueco
implementé \texttt{scripts/generate\_multi\_t.py}, que genera 11 capas
intermedias ($t = 0.0, 0.1, \ldots, 1.0$) con la coordenada $z$ real
interpolada linealmente entre los dos slices originales:
$z(t) = z_A + t \cdot (z_B - z_A)$.

\begin{lstlisting}[language=bash,basicstyle=\small\ttfamily]
# Genera 11 contornos interpolados con z real entre los slices 70 y 72
python scripts/generate_multi_t.py \
    data/contours/BraTS-GLI-00008-100/slice_0070.obj \
    data/contours/BraTS-GLI-00008-100/slice_0072.obj \
    70 72
\end{lstlisting}

Cada corrida deposita los 11 \texttt{.obj} más un
\texttt{preview\_3d.png} y un \texttt{index.csv} (columnas
\texttt{t, z, filename, n\_vertices, self\_int}) en una subcarpeta con
\emph{timestamp} dentro de \texttt{src/interpolation-lineal/output/multi\_t/}.
Este \texttt{index.csv} es el punto de integración con el pipeline de
reconstrucción de Santiago: él consume directamente el listado para apilar
los contornos y alimentar su superficie de marching cubes.
